\subsection{Two binary adjacent stalks and the actual long-edge attachment}

\begin{theorem}[Proved: opposite three-element bands with binary adjacent stalks]
\label{CP2:A-binary-adjacent-opposite-three-bands}
Let $L_3=\{1,e_0,e_1\}$ with $e_i e_j=e_i$, and let
$R_3=\{1,f_0,f_1\}$ with $f_i f_j=f_j$.
Suppose a category has full linear support $0<1<2$, corners
$L_3,1,R_3$, two arrows in each adjacent hom-set, and any
nonempty finite long hom-set. All rational degree-two products
of its classifying space vanish.
\end{theorem}
\begin{proof}
Write $S=C(0,1)$, $T=C(1,2)$, $U=C(0,2)$, and
$\mu:T\times S\to U$ for composition. The right $L_3$ actors
on $S$ are two retractions with the same image, and the same is
true of the left $R_3$ actors on $T$. On a two-point set they
are therefore either both identities or the same constant map.
Indeed the two possible distinct constants do not satisfy
$PQ=Q$, $QP=P$, whereas the identity cannot occur with a
constant in such a pair.

If the actors on $S$ are identities, its element space is the
disjoint union of two copies of $BL_3$. Each copy is contractible,
since $L_3$ has a left zero. The full $(0,1)$ edge consequently
has zero rational second homology. The opposite corner $BR_3$
is also contractible, so Lemma~\ref{CP2:A-two-block} kills
the products. Treat identity actors on $T$ dually. It remains
to consider the equal-constant actions. Label the fixed states
$s_0,t_0$, and the other states $s_1,t_1$.

Put $R=\mathbb QL_3$, $B=\mathbb QR_3$,
$a=1-e_0$, and $b=1-f_0$. The following are exact projective
resolutions, on the displayed sides:
\[
\begin{split}
 0\longrightarrow Ra\longrightarrow Re_0
   \longrightarrow\mathbb Q\longrightarrow0,
 &\qquad a\longmapsto e_1-e_0,\\
 0\longrightarrow bB\longrightarrow f_0B
   \longrightarrow\mathbb Q\longrightarrow0,
 &\qquad b\longmapsto f_1-f_0.
\end{split}
\]
The ideals in degree one have dimension one and those in
degree zero have dimension two. Their displayed images are
exactly the augmentation kernels. The basis elements $a,b$
are idempotents, so all terms are projective. These formulas
also make every subsequent comparison natural in a biset.

Let $A=B(S,L_3)$, $D=B(R_3,T)$, and
$K=B(R_3,U,L_3)$ be the actual edge bar spaces. The first
resolution computes $A$ by the two-term complex
\[
 \mathbb Q\{s_1-s_0\}[1]\oplus\mathbb Q\{s_0\}[0]
\]
with zero differential; the second gives the identical
description of $D$ using $t_0,t_1$. In particular each space
is connected, has one-dimensional rational first homology,
and has no higher rational homology. This does not assert
that the constant two-state action space is acyclic.

Tensor the two projective resolutions with $\mathbb QU$.
It computes the rational homology of $K$ by a natural complex
concentrated in degrees zero, one and two, whose degree-two term is
\[
 C_2(K)=b\,\mathbb QU\,a.
\]
In particular $H_j(K;\mathbb Q)=0$ for $j\ge3$, and
$H_2(K;\mathbb Q)$ is the kernel of its degree-two differential,
as an actual subspace of $\mathbb QU$. There is no incoming
degree-three boundary in this projective complex.

Since the middle corner is trivial, the actual triangular
bar face space is $F=D\times A$. Its maps to the two adjacent
edge spaces are the projections, and its map to $K$ is induced
by the actual composition $\mu$. Set $u_{ij}=\mu(t_i,s_j)$.
The natural comparison of the displayed resolutions sends
the generator of $H_2(F;\mathbb Q)$ to the cycle
\[
 v=u_{11}-u_{10}-u_{01}+u_{00}\ \in b\,\mathbb QU\,a.
\]
To see the formula directly, tensor
$(t_1-t_0)\otimes(s_1-s_0)$ and apply $\mu$ linearly.
Equivariance gives $f_0u_{ij}=f_1u_{ij}=u_{0j}$ and
$u_{ij}e_0=u_{ij}e_1=u_{i0}$, so the vector belongs to the
displayed term and both components of its differential vanish.

Filter the compatible triangular bar model of $BC$ by its
number of strict object changes. The filtration has columns
zero, one and two: the three corner spaces, the edge spaces
$A,D,K$, and $F$. All corner spaces have zero positive
rational homology. The preceding calculations therefore give
\[
 H_4(BC;\mathbb Q)
   \cong\ker\bigl(H_2(F;\mathbb Q)
                    \xrightarrow{\mu_*}H_2(K;\mathbb Q)\bigr).
\]
There is no other total-degree-four term and no possible
nonzero higher filtration differential at that position.
If $v\ne0$, the source is one-dimensional, this map is
injective, and $H^4(BC;\mathbb Q)=0$.

Suppose instead that $v=0$. As the states of $U$ form a free
basis, this says that the multisets
$\{u_{11},u_{00}\}$ and $\{u_{10},u_{01}\}$ agree.
Either $u_{00}=u_{01}$ and $u_{10}=u_{11}$, or
$u_{00}=u_{10}$ and $u_{01}=u_{11}$.
Thus $\mu$ ignores one of its two binary coordinates.
Assume it ignores $S$; the other case is dual. The map
$H_1(A;\mathbb Q)\to H_1(K;\mathbb Q)$ obtained by
fixing $t_0$ sends its generator to $u_{01}-u_{00}=0$.

For completeness, the same actual filtration in degree two
has the edge term
$H_1(A)\oplus H_1(D)\oplus H_1(K)$ modulo the image of
$H_1(F)$. The filtration term $H_0(F)$ in total degree two
has zero kernel: its boundary has nonzero components in
both $H_0(A)$ and $H_0(D)$. The $A$-summand of $H_1(F)$
maps identically to $H_1(A)$, trivially to $H_1(D)$, and
by the zero map just computed to $H_1(K)$.
It follows, by duality and the natural edge inclusions, that
every class in $H^2(BC;\mathbb Q)$ restricts to zero on
the full $(0,1)$ edge. Every such class also restricts to
zero on the corner $BR_3$. The relative-product proof of
Lemma~\ref{CP2:A-two-block} then kills the product of any
two global second classes. This proves the assertion in
the remaining case as well.
\end{proof}

\begin{corollary}[Proved: the small-endpoint binary-adjacent frontier at nineteen]
\label{CP2:A-small-endpoints-binary-adjacent-nineteen}
A nineteen-arrow full-linear three-object category with
trivial middle corner, both other corners of order at most
three, and both adjacent hom-sets of size two cannot have
type $\CP^2$. In particular the profile
$(3,1,3;2,2,8)$ is excluded.
\end{corollary}
\begin{proof}
All corners of order at most three have zero positive
rational cohomology. If an endpoint is a group, its corner
algebra and that of the middle object are semisimple; their
full edge has zero second homology, and the opposite corner
has zero second homology. The two-block argument applies.
For a nongroup source whose trivial left module is projective,
or a nongroup sink whose trivial right module is projective,
use Corollary~\ref{CP2:A-extreme-minimal-group-nineteen}.
The complete monoid classification through order three
leaves only $L_3$ at the source and $R_3$ at the sink.
Their remaining case is the preceding theorem.
\end{proof}

The separate literal check binds the complete order-at-most-three
catalogs, the two projective resolutions, every binary band action,
and all partitions of the four composition values. The earlier
failed assertion that a constant two-state $R_3$ action has zero
first homology remains withdrawn; the present proof explicitly
uses its nonzero first homology and the actual triangular attachment.
